\begin{figure*}[ht]
    \centering
    %\includegraphics[width=0.95\textwidth]{figures/ch3_aoc/illustrations/fig_method_2.png}
    %\includegraphics[width=0.95\textwidth]{figures/ch3_aoc/illustrations/fig_method_3b.png}
    \includegraphics[width=0.95\textwidth]{figures/ch3_aoc/illustrations/method_v4.png}
    \caption{\textbf{AOC enables the construction of orthogonal kernels with customizable sizes and strides.} By leveraging the $\mathbconv$ operator, we can fuse kernels obtained from two existing methods, namely BCOP and RKO. With this approach, we can build orthogonal kernels that support native striding, effectively mitigating the drawbacks of the two base methods. %\thib{Minor: Moreover, the caption is not sufficiently informative to help readers understand the figure.}
    }% \vspace{-4mm}}
    \label{fig:method_bigpic}
\end{figure*}
